%!TeX root = ../main.tex

% ==================================================================================================================================
% 7. Resultados e discussão
% ==================================================================================================================================
\section{Resultados e discussão}

Este capítulo apresenta e discute os resultados obtidos com o ambiente de testes
descrito no capítulo anterior. Cada problema é analisado quanto ao tempo de
execução e, quando há heurísticas, quanto à qualidade da solução. As tabelas de
tempo separam as implementações em linguagem interpretada das de núcleo
compilado, conforme justificado na Seção~\ref{sec:criterios}, e todos os valores
são gerados diretamente a partir dos dados medidos. Ao final, discutem-se os
defeitos encontrados em implementações públicas e o crescimento empírico dos
tempos frente à complexidade teórica.

% ---------------------------------------------------------------------------
\subsection{Cardinalidade máxima em grafos bipartidos}

As Tabelas~\ref{tab:mcm-grafo-3-denso-com-perfeito-mcm-n-tempo-python}
e~\ref{tab:mcm-grafo-3-denso-com-perfeito-mcm-n-tempo-c} apresentam os tempos dos
algoritmos para o grafo denso, separados por linguagem de implementação, e a
Tabela~\ref{tab:mcm-grafo-3-denso-com-perfeito-mcm-n-qualidade}, a qualidade das
soluções.

\begin{table}[H]
    \centering
    \caption{Tempo de execução em Python — Grafo 3 — Denso   \textbar{} COM perfeito \textbar{} MCM = n (mediana, ms)}
    \label{tab:mcm-grafo-3-denso-com-perfeito-mcm-n-tempo-python}
    \renewcommand{\arraystretch}{1.2}
    \begin{tabular}{@{} l l l l rrrrr @{}}
        \toprule
        \textbf{Algoritmo} & \textbf{Tipo} & \textbf{Tempo} & \textbf{Espaço} & \textbf{$n=50$} & \textbf{$n=100$} & \textbf{$n=150$} & \textbf{$n=200$} & \textbf{$n=250$} \\
        \midrule
        Aug. Paths (wbchristerson) & exato & $O(V \cdot E)$ & $O(V+E)$ & 2.08 & 19.12 & 34.10 & 39.79 & 259.82 \\
        Hopcroft-Karp (sofiat) & exato & $O(E \cdot \sqrt{V})$ & $O(V+E)$ & 1.14 & 2.41 & 3.99 & 6.19 & 7.53 \\
        Edmonds-Karp (Maxflow-Algs) & exato & $O(V \cdot E^{2})$ & $O(V^{2})$ & 12.19 & 75.08 & 238.34 & 556.37 & 1077.90 \\
        Hopcroft-Karp (NetworkX) & baseline & $O(E \cdot \sqrt{V})$ & $O(V+E)$ & 1.51 & 3.15 & 4.84 & 6.94 & 8.66 \\
        Gulosa & heuristica & $O(E)$ & $O(V)$ & 0.03 & 0.06 & 0.08 & 0.11 & 0.14 \\
        Karp-Sipser & heuristica & $O(V \cdot E)$ & $O(V+E)$ & 0.49 & 1.35 & 3.87 & 4.65 & 6.86 \\
        \bottomrule
    \end{tabular}
\end{table}

\begin{table}[ht]
    \centering
    \caption{Tempo de execução em compilado (C/C++) — Grafo 3 — Denso   \textbar{} COM perfeito \textbar{} MCM = n (mediana, ms)}
    \label{tab:mcm-grafo-3-denso-com-perfeito-mcm-n-tempo-c}
    \renewcommand{\arraystretch}{1.2}
    \begin{tabular}{@{} l l l l rrrrr @{}}
        \toprule
        \textbf{Algoritmo} & \textbf{Tipo} & \textbf{Tempo} & \textbf{Espaço} & \textbf{$n=50$} & \textbf{$n=100$} & \textbf{$n=150$} & \textbf{$n=200$} & \textbf{$n=250$} \\
        \midrule
        Bipartite Match. (igraph) & baseline & $O(E \cdot \sqrt{V})$ & $O(V+E)$ & 0.20 & 0.36 & 0.61 & 0.85 & 1.01 \\
        \bottomrule
    \end{tabular}
\end{table}

\begin{table}[H]
    \centering
    \caption{Qualidade da solução — Grafo 3 — Denso   \textbar{} COM perfeito \textbar{} MCM = n}
    \label{tab:mcm-grafo-3-denso-com-perfeito-mcm-n-qualidade}
    \renewcommand{\arraystretch}{1.2}
    \begin{tabular}{@{} l l l rrrrr @{}}
        \toprule
        \textbf{Algoritmo} & \textbf{Tipo} & \textbf{Tempo} & \textbf{$n=50$} & \textbf{$n=100$} & \textbf{$n=150$} & \textbf{$n=200$} & \textbf{$n=250$} \\
        \midrule
        Aug. Paths (wbchristerson) & exato & $O(V \cdot E)$ & 50 & 100 & 150 & 200 & 250 \\
        Hopcroft-Karp (sofiat) & exato & $O(E \cdot \sqrt{V})$ & 50 & 100 & 150 & 200 & 250 \\
        Edmonds-Karp (Maxflow-Algs) & exato & $O(V \cdot E^{2})$ & 50 & 100 & 150 & 200 & 250 \\
        Hopcroft-Karp (NetworkX) & baseline & $O(E \cdot \sqrt{V})$ & 50 & 100 & 150 & 200 & 250 \\
        Bipartite Match. (igraph) & baseline & $O(E \cdot \sqrt{V})$ & 50 & 100 & 150 & 200 & 250 \\
        Gulosa & heuristica & $O(E)$ & 47 & 93 & 142 & 191 & 235 \\
        Karp-Sipser & heuristica & $O(V \cdot E)$ & 50 & 100 & 150 & 200 & 249 \\
        \bottomrule
    \end{tabular}
\end{table}


Entre os algoritmos exatos, a diferença de desempenho é de mais de duas ordens de
grandeza: o Edmonds-Karp, que reduz o emparelhamento a um problema de fluxo e
reconstrói a busca a cada aumento, é dramaticamente mais lento que o
Hopcroft-Karp, que agrupa os caminhos aumentantes em fases. Essa distância
confirma na prática o que a complexidade assintótica prevê e ilustra por que, em
grafos maiores, a escolha do algoritmo importa muito mais do que a constante de
implementação. A implementação de núcleo compilado (igraph), na tabela dos
compilados, é a mais rápida de todas, como esperado.

Quanto às heurísticas, os resultados evidenciam o compromisso central entre
qualidade e tempo. A heurística gulosa é a mais rápida de todo o conjunto, porém
não atinge o ótimo em todas as instâncias. A heurística de Karp-Sipser, ao custo
de um tempo comparável ao dos exatos mais eficientes, encontra o ótimo em todas
as instâncias avaliadas --- reflexo da regra dos vértices de grau um, que fixa
escolhas comprovadamente ótimas e, em grafos esparsos, chega a dispensar o passo
guloso.

% ---------------------------------------------------------------------------
\subsection{Atribuição linear}

Os tempos dos algoritmos para a matriz densa constam na
Tabela~\ref{tab:assignment-matriz-1-densa-uniforme-custo-in-1-100-tempo-python}
(interpretados) e na
Tabela~\ref{tab:assignment-matriz-1-densa-uniforme-custo-in-1-100-tempo-c}
(compilados); a Tabela~\ref{tab:assignment-matriz-1-densa-uniforme-custo-in-1-100-qualidade}
traz a qualidade, na qual se veem também as células em que dois dos algoritmos
falham.

\begin{table}[H]
    \centering
    \caption{Tempo de execução em Python — Matriz 1 — Densa uniforme  \textbar{} custo in [1,100] (mediana, ms)}
    \label{tab:assignment-matriz-1-densa-uniforme-custo-in-1-100-tempo-python}
    \renewcommand{\arraystretch}{1.2}
    \begin{tabular}{@{} l l l l rrrrrrrrr @{}}
        \toprule
        \textbf{Algoritmo} & \textbf{Tipo} & \textbf{Tempo} & \textbf{Espaço} & \textbf{$n=10$} & \textbf{$n=20$} & \textbf{$n=30$} & \textbf{$n=50$} & \textbf{$n=75$} & \textbf{$n=100$} & \textbf{$n=150$} & \textbf{$n=200$} & \textbf{$n=250$} \\
        \midrule
        Húngaro (munkres) & exato & $O(n^{3})$ & $O(n^{2})$ & 0.39 & 1.98 & 4.82 & 15.43 & 41.61 & 97.23 & 239.96 & 428.17 & 679.84 \\
        Húngaro (benchaplin) & exato & $O(n^{3})$ & $O(n^{2})$ & 11.11 & 92.32 & 391.64 & BugAttr & BugAttr & BugAttr & BugAttr & BugAttr & BugAttr \\
        Húngaro (hunalgorithm) & exato & $O(n^{3})$ & $O(n^{2})$ & 2.10 & BugPerc & BugPerc & BugPerc & BugPerc & BugPerc & BugPerc & BugPerc & BugPerc \\
        Húngaro (mayorx/KM) & exato & $O(n^{3})$ & $O(n^{2})$ & 0.99 & 2.25 & 7.49 & 12.30 & 21.11 & 29.30 & 58.44 & 58.55 & 86.71 \\
        SSP (flows) & exato & $O(n^{4})$* & $O(n^{2})$ & 0.45 & 2.65 & 6.92 & 29.97 & 98.01 & 235.13 & — & — & — \\
        \bottomrule
    \end{tabular}
\end{table}

\begin{table}[H]
    \centering
    \caption{Tempo de execução em compilado (C/C++) — Matriz 1 — Densa uniforme  \textbar{} custo in [1,100] (mediana, ms)}
    \label{tab:assignment-matriz-1-densa-uniforme-custo-in-1-100-tempo-c}
    \renewcommand{\arraystretch}{1.2}
    \footnotesize
    \adjustbox{max width=\textwidth}{%
    \begin{tabular}{@{} l l l l rrrrrrrrr @{}}
        \toprule
        \textbf{Algoritmo} & \textbf{Tipo} & \textbf{Tempo} & \textbf{Espaço} & \textbf{$n=10$} & \textbf{$n=20$} & \textbf{$n=30$} & \textbf{$n=50$} & \textbf{$n=75$} & \textbf{$n=100$} & \textbf{$n=150$} & \textbf{$n=200$} & \textbf{$n=250$} \\
        \midrule
        SciPy (linear\_sum\_assg) & baseline & $O(n^{3})$ & $O(n^{2})$ & 0.02 & 0.04 & 0.08 & 0.23 & 0.51 & 0.84 & 2.19 & 3.68 & 5.48 \\
        Jonker-Volgenant (lap) & baseline & $O(n^{3})$ & $O(n^{2})$ & 0.02 & 0.05 & 0.07 & 0.17 & 0.36 & 0.61 & 1.36 & 2.31 & 3.73 \\
        \bottomrule
    \end{tabular}%
    }
\end{table}

\begin{table}[ht]
    \centering
    \caption{Qualidade da solução — Matriz 1 — Densa uniforme  \textbar{} custo in [1,100]}
    \label{tab:assignment-matriz-1-densa-uniforme-custo-in-1-100-qualidade}
    \renewcommand{\arraystretch}{1.2}
    \begin{tabular}{@{} l l l rrrrrrrrr @{}}
        \toprule
        \textbf{Algoritmo} & \textbf{Tipo} & \textbf{Tempo} & \textbf{$n=10$} & \textbf{$n=20$} & \textbf{$n=30$} & \textbf{$n=50$} & \textbf{$n=75$} & \textbf{$n=100$} & \textbf{$n=150$} & \textbf{$n=200$} & \textbf{$n=250$} \\
        \midrule
        Húngaro (munkres) & exato & $O(n^{3})$ & 102 & 146 & 168 & 217 & 238 & 244 & 258 & 275 & 300 \\
        Húngaro (benchaplin) & exato & $O(n^{3})$ & 102 & 146 & 168 & BugAttr & BugAttr & BugAttr & BugAttr & BugAttr & BugAttr \\
        Húngaro (hunalgorithm) & exato & $O(n^{3})$ & 102 & BugPerc & BugPerc & BugPerc & BugPerc & BugPerc & BugPerc & BugPerc & BugPerc \\
        Húngaro (mayorx/KM) & exato & $O(n^{3})$ & 102 & 146 & 168 & 217 & 238 & 244 & 258 & 275 & 300 \\
        SSP (flows) & exato & $O(n^{4})$* & 102 & 146 & 168 & 217 & 238 & 244 & — & — & — \\
        SciPy (linear\_sum\_assg) & baseline & $O(n^{3})$ & 102 & 146 & 168 & 217 & 238 & 244 & 258 & 275 & 300 \\
        Jonker-Volgenant (lap) & baseline & $O(n^{3})$ & 102 & 146 & 168 & 217 & 238 & 244 & 258 & 275 & 300 \\
        \bottomrule
    \end{tabular}
\end{table}


Todos os algoritmos que concluem atingem o custo ótimo, de modo que a comparação
recai sobre o tempo. Entre as quatro implementações do método Húngaro, obtidas de
autores distintos para exatamente o mesmo algoritmo, observa-se uma variação de
mais de uma ordem de grandeza --- evidência de que a qualidade da implementação,
e não apenas o algoritmo, determina o desempenho na prática. As implementações de
núcleo compilado (SciPy e Jonker-Volgenant) são, novamente, as mais rápidas.

Duas das implementações do Húngaro apresentam defeitos que impedem sua conclusão
em parte das instâncias densas; como esse é um resultado relevante por si só, ele
é tratado em conjunto na Seção~\ref{sec:defeitos}.

% ---------------------------------------------------------------------------
\subsection{Atribuição gargalo}

A Tabela~\ref{tab:bottleneck-matriz-1-densa-uniforme-custo-in-1-100-tempo-python}
compara os tempos dos dois algoritmos exatos para a matriz densa, e a
Tabela~\ref{tab:bottleneck-matriz-1-densa-uniforme-custo-in-1-100-qualidade}, o
custo gargalo que encontram.

\begin{table}[H]
    \centering
    \caption{Tempo de execução em Python — Matriz 1 — Densa uniforme  \textbar{} custo in [1,100] (mediana, ms)}
    \label{tab:bottleneck-matriz-1-densa-uniforme-custo-in-1-100-tempo-python}
    \renewcommand{\arraystretch}{1.2}
    \footnotesize
    \adjustbox{max width=\textwidth}{%
    \begin{tabular}{@{} l l l l rrrrrrrr @{}}
        \toprule
        \textbf{Algoritmo} & \textbf{Tipo} & \textbf{Tempo} & \textbf{Espaço} & \textbf{$n=10$} & \textbf{$n=20$} & \textbf{$n=30$} & \textbf{$n=50$} & \textbf{$n=75$} & \textbf{$n=100$} & \textbf{$n=150$} & \textbf{$n=200$} \\
        \midrule
        Threshold (busca binária) & exato & $O(n^{3} log n)$ & $O(n^{2})$ & 0.10 & 0.50 & 1.44 & 4.48 & 13.50 & 32.06 & 109.54 & 257.24 \\
        Caminhos aumentantes & exato & $O(n^{3})$ & $O(n^{2})$ & 0.05 & 0.23 & 0.62 & 2.10 & 5.43 & 11.32 & 32.65 & 68.21 \\
        \bottomrule
    \end{tabular}%
    }
\end{table}

\begin{table}[ht]
    \centering
    \caption{Qualidade da solução — Matriz 1 — Densa uniforme  \textbar{} custo in [1,100]}
    \label{tab:bottleneck-matriz-1-densa-uniforme-custo-in-1-100-qualidade}
    \renewcommand{\arraystretch}{1.2}
    \begin{tabular}{@{} l l l rrrrrrrr @{}}
        \toprule
        \textbf{Algoritmo} & \textbf{Tipo} & \textbf{Tempo} & \textbf{$n=10$} & \textbf{$n=20$} & \textbf{$n=30$} & \textbf{$n=50$} & \textbf{$n=75$} & \textbf{$n=100$} & \textbf{$n=150$} & \textbf{$n=200$} \\
        \midrule
        Threshold (busca binária) & exato & $O(n^{3} log n)$ & 26 & 18 & 18 & 13 & 8 & 7 & 5 & 4 \\
        Caminhos aumentantes & exato & $O(n^{3})$ & 26 & 18 & 18 & 13 & 8 & 7 & 5 & 4 \\
        \bottomrule
    \end{tabular}
\end{table}


Ambos produzem sempre o mesmo custo gargalo ótimo, o que serve também como
verificação cruzada de corretude. O método de caminhos aumentantes, construtivo,
mostra-se consistentemente mais rápido que o método de \textit{threshold}, que
repete uma verificação de emparelhamento perfeito a cada passo da busca binária
sobre o limiar. A vantagem do método construtivo cresce com o tamanho da
instância.

% ---------------------------------------------------------------------------
\subsection{Casamento estável}

A Tabela~\ref{tab:stable-marriage-perfil-1-aleatorio-permutacoes-uniformes-tempo-python}
apresenta os tempos das três variantes de Gale-Shapley para preferências
aleatórias, e a
Tabela~\ref{tab:stable-marriage-perfil-1-aleatorio-permutacoes-uniformes-qualidade},
a verificação de estabilidade e de concordância entre elas.

\begin{table}[H]
    \centering
    \caption{Tempo de execução em Python — Perfil 1 — Aleatório       \textbar{} permutações uniformes (mediana, ms)}
    \label{tab:stable-marriage-perfil-1-aleatorio-permutacoes-uniformes-tempo-python}
    \renewcommand{\arraystretch}{1.2}
    \begin{tabular}{@{} l l l l rrrrrrr @{}}
        \toprule
        \textbf{Algoritmo} & \textbf{Tipo} & \textbf{Tempo} & \textbf{Espaço} & \textbf{$n=25$} & \textbf{$n=50$} & \textbf{$n=75$} & \textbf{$n=100$} & \textbf{$n=200$} & \textbf{$n=300$} & \textbf{$n=400$} \\
        \midrule
        Gale-Shapley (OOP/pip) & exato & $O(n^{2})$† & $O(n^{2})$ & 7.54 & 94.14 & 321.96 & 692.09 & — & — & — \\
        Gale-Shapley (numérico/pip) & exato & $O(n^{2})$ & $O(n^{2})$ & 0.44 & 1.57 & 3.36 & 5.52 & 21.70 & 50.59 & 95.89 \\
        Gale-Shapley (Lattas/dict) & exato & $O(n^{3})$* & $O(n^{2})$ & 0.31 & 1.29 & 2.85 & 5.08 & 20.07 & 52.20 & 100.25 \\
        \bottomrule
    \end{tabular}
\end{table}

\begin{table}[ht]
    \centering
    \caption{Qualidade da solução — Perfil 1 — Aleatório       \textbar{} permutações uniformes}
    \label{tab:stable-marriage-perfil-1-aleatorio-permutacoes-uniformes-qualidade}
    \renewcommand{\arraystretch}{1.2}
    \begin{tabular}{@{} l l l rrrrrrr @{}}
        \toprule
        \textbf{Algoritmo} & \textbf{Tipo} & \textbf{Tempo} & \textbf{$n=25$} & \textbf{$n=50$} & \textbf{$n=75$} & \textbf{$n=100$} & \textbf{$n=200$} & \textbf{$n=300$} & \textbf{$n=400$} \\
        \midrule
        Gale-Shapley (OOP/pip) & exato & $O(n^{2})$† & 0 =ref & 0 =ref & 0 =ref & 0 =ref & — & — & — \\
        Gale-Shapley (numérico/pip) & exato & $O(n^{2})$ & 0 =ref & 0 =ref & 0 =ref & 0 =ref & 0 =ref & 0 =ref & 0 =ref \\
        Gale-Shapley (Lattas/dict) & exato & $O(n^{3})$* & 0 =ref & 0 =ref & 0 =ref & 0 =ref & 0 =ref & 0 =ref & 0 =ref \\
        \bottomrule
    \end{tabular}
\end{table}


Pela unicidade do emparelhamento estável ótimo para os pretendentes, as três
variantes produzem exatamente o mesmo resultado; a qualidade, portanto, não as
distingue. A variante orientada a objetos, embora teoricamente quadrática,
mostra-se substancialmente mais lenta e inviável nos tamanhos maiores, por conta
do sobrecusto de reconstruir estruturas a cada rodada --- comportamento
quantificado adiante, na análise de crescimento empírico. As variantes que operam
sobre vetores e dicionários escalam sem dificuldade até os maiores tamanhos
avaliados.

% ---------------------------------------------------------------------------
\subsection{Emparelhamento em grafos gerais}

Para cardinalidade máxima, comparam-se o algoritmo de Blossom e a biblioteca
NetworkX
(Tabelas~\ref{tab:general-cardinalidade-esparso-grau-3-tempo-python}
e~\ref{tab:general-cardinalidade-esparso-grau-3-qualidade}); para peso máximo, a
implementação de van Rantwijk, a mesma biblioteca e a heurística Path Growing
(Tabelas~\ref{tab:general-peso-esparso-grau-3-tempo-python}
e~\ref{tab:general-peso-esparso-grau-3-qualidade}).

\begin{table}[H]
    \centering
    \caption{Tempo de execução em Python — Cardinalidade — esparso (grau \textasciitilde{}3) (mediana, ms)}
    \label{tab:general-cardinalidade-esparso-grau-3-tempo-python}
    \renewcommand{\arraystretch}{1.2}
    \footnotesize
    \adjustbox{max width=\textwidth}{%
    \begin{tabular}{@{} l l l l rrrrrr @{}}
        \toprule
        \textbf{Algoritmo} & \textbf{Tipo} & \textbf{Tempo} & \textbf{Espaço} & \textbf{$n=10$} & \textbf{$n=20$} & \textbf{$n=30$} & \textbf{$n=50$} & \textbf{$n=75$} & \textbf{$n=100$} \\
        \midrule
        Blossom (Edmonds) & exato & $O(V^{3})$ & $O(V^{2})$ & 1.66 & 3.21 & 3.97 & 20.36 & 43.90 & 103.46 \\
        NetworkX (cardinalidade) & baseline & $O(V^{3})$ & $O(V+E)$ & 0.40 & 0.63 & 0.74 & 1.41 & 2.28 & 3.15 \\
        \bottomrule
    \end{tabular}%
    }
\end{table}

\begin{table}[ht]
    \centering
    \caption{Qualidade da solução — Cardinalidade — esparso (grau \textasciitilde{}3)}
    \label{tab:general-cardinalidade-esparso-grau-3-qualidade}
    \renewcommand{\arraystretch}{1.2}
    \begin{tabular}{@{} l l l rrrrrr @{}}
        \toprule
        \textbf{Algoritmo} & \textbf{Tipo} & \textbf{Tempo} & \textbf{$n=10$} & \textbf{$n=20$} & \textbf{$n=30$} & \textbf{$n=50$} & \textbf{$n=75$} & \textbf{$n=100$} \\
        \midrule
        Blossom (Edmonds) & exato & $O(V^{3})$ & 5 & 9 & 13 & 23 & 31 & 43 \\
        NetworkX (cardinalidade) & baseline & $O(V^{3})$ & 5 & 9 & 13 & 23 & 31 & 43 \\
        \bottomrule
    \end{tabular}
\end{table}

\begin{table}[H]
    \centering
    \caption{Tempo de execução em Python — Peso — esparso (grau \textasciitilde{}3) (mediana, ms)}
    \label{tab:general-peso-esparso-grau-3-tempo-python}
    \renewcommand{\arraystretch}{1.2}
    \begin{tabular}{@{} l l l l rrrrrr @{}}
        \toprule
        \textbf{Algoritmo} & \textbf{Tipo} & \textbf{Tempo} & \textbf{Espaço} & \textbf{$n=10$} & \textbf{$n=20$} & \textbf{$n=30$} & \textbf{$n=50$} & \textbf{$n=75$} & \textbf{$n=100$} \\
        \midrule
        van Rantwijk (mwmatching) & exato & $O(V^{3})$ & $O(V+E)$ & 0.26 & 0.54 & 0.86 & 2.19 & 3.49 & 6.30 \\
        NetworkX (peso máximo) & baseline & $O(V^{3})$ & $O(V+E)$ & 0.59 & 1.25 & 1.98 & 4.87 & 7.52 & 13.58 \\
        Path Growing & heuristica & $O(E)$ & $O(V+E)$ & 0.03 & 0.05 & 0.07 & 0.11 & 0.17 & 0.23 \\
        \bottomrule
    \end{tabular}
\end{table}

\begin{table}[H]
    \centering
    \caption{Qualidade da solução — Peso — esparso (grau \textasciitilde{}3)}
    \label{tab:general-peso-esparso-grau-3-qualidade}
    \renewcommand{\arraystretch}{1.2}
    \begin{tabular}{@{} l l l rrrrrr @{}}
        \toprule
        \textbf{Algoritmo} & \textbf{Tipo} & \textbf{Tempo} & \textbf{$n=10$} & \textbf{$n=20$} & \textbf{$n=30$} & \textbf{$n=50$} & \textbf{$n=75$} & \textbf{$n=100$} \\
        \midrule
        van Rantwijk (mwmatching) & exato & $O(V^{3})$ & 309 & 604 & 790 & 1284 & 1907 & 2929 \\
        NetworkX (peso máximo) & baseline & $O(V^{3})$ & 309 & 604 & 790 & 1284 & 1907 & 2929 \\
        Path Growing & heuristica & $O(E)$ & 278 & 510 & 732 & 1030 & 1492 & 2407 \\
        \bottomrule
    \end{tabular}
\end{table}


No caso ponderado, um resultado é digno de nota: a implementação independente de
van Rantwijk é mais rápida que a rotina correspondente do NetworkX, ainda que
esta última seja uma adaptação daquele mesmo código. A comparação, portanto, não
opõe duas abordagens, e sim uma implementação a um descendente empacotado dela;
a diferença observada mede o efeito das camadas de generalidade que a biblioteca
acrescenta. A heurística Path Growing, por sua vez, é ordens de grandeza mais
rápida que os exatos e entrega uma fração elevada do peso ótimo, bem acima da sua
garantia teórica de metade --- mais um exemplo de que, na prática, a heurística
rende muito mais do que o seu pior caso.

% ---------------------------------------------------------------------------
\subsection{Emparelhamento induzido máximo}

Sendo NP-difícil, o emparelhamento induzido máximo tem um perfil distinto dos
demais problemas: o algoritmo exato só é viável em instâncias pequenas, e a
qualidade das heurísticas passa a ser o eixo central da comparação. A
Tabela~\ref{tab:mim-grafo-esparso-grau-3-tempo-python} traz os tempos e a
Tabela~\ref{tab:mim-grafo-esparso-grau-3-qualidade}, a qualidade.

\begin{table}[H]
    \centering
    \caption{Tempo de execução em Python — Grafo esparso (grau \textasciitilde{}3) (mediana, ms)}
    \label{tab:mim-grafo-esparso-grau-3-tempo-python}
    \renewcommand{\arraystretch}{1.2}
    \footnotesize
    \adjustbox{max width=\textwidth}{%
    \begin{tabular}{@{} l l l l rrrrrrrrrr @{}}
        \toprule
        \textbf{Algoritmo} & \textbf{Tipo} & \textbf{Tempo} & \textbf{Espaço} & \textbf{$n=10$} & \textbf{$n=15$} & \textbf{$n=20$} & \textbf{$n=25$} & \textbf{$n=30$} & \textbf{$n=40$} & \textbf{$n=50$} & \textbf{$n=100$} & \textbf{$n=200$} & \textbf{$n=400$} \\
        \midrule
        Exato (MIS conflitos) & exato & $exp.$ & $O(E^{2})$ & 0.16 & 0.29 & 1.49 & 1.99 & 5.10 & 74.61 & — & — & — & — \\
        Guloso aleatório & heuristica & $O(E^{2})$ & $O(V+E)$ & 0.05 & 0.05 & 0.08 & 0.13 & 0.18 & 0.29 & 0.47 & 1.50 & 8.28 & 27.29 \\
        Guloso grau mínimo & heuristica & $O(E^{2})$ & $O(E^{2})$ & 0.18 & 0.36 & 0.78 & 1.47 & 1.83 & 2.78 & 4.27 & 18.48 & 85.90 & 378.36 \\
        \bottomrule
    \end{tabular}%
    }
\end{table}

\begin{table}[H]
    \centering
    \caption{Qualidade da solução — Grafo esparso (grau \textasciitilde{}3)}
    \label{tab:mim-grafo-esparso-grau-3-qualidade}
    \renewcommand{\arraystretch}{1.2}
    \footnotesize
    \adjustbox{max width=\textwidth}{%
    \begin{tabular}{@{} l l l rrrrrrrrrr @{}}
        \toprule
        \textbf{Algoritmo} & \textbf{Tipo} & \textbf{Tempo} & \textbf{$n=10$} & \textbf{$n=15$} & \textbf{$n=20$} & \textbf{$n=25$} & \textbf{$n=30$} & \textbf{$n=40$} & \textbf{$n=50$} & \textbf{$n=100$} & \textbf{$n=200$} & \textbf{$n=400$} \\
        \midrule
        Exato (MIS conflitos) & exato & $exp.$ & 2 & 3 & 5 & 5 & 8 & 10 & — & — & — & — \\
        Guloso aleatório & heuristica & $O(E^{2})$ & 2 & 2 & 3 & 3 & 5 & 7 & 10 & 18 & 41 & 75 \\
        Guloso grau mínimo & heuristica & $O(E^{2})$ & 2 & 3 & 5 & 5 & 8 & 10 & 11 & 25 & 49 & 93 \\
        \bottomrule
    \end{tabular}%
    }
\end{table}


O algoritmo exato, que reduz o problema a conjunto independente máximo no grafo de
conflitos, resolve instâncias pequenas em poucos milissegundos, mas seu custo
cresce rapidamente e o torna inviável além de algumas dezenas de vértices --- o
que justifica empiricamente a necessidade das heurísticas. Entre estas, evidencia-se
um compromisso claro: a heurística de grau mínimo encontra soluções de qualidade
consistentemente superior à da escolha aleatória, ao custo de um tempo bem maior,
por manter e atualizar o grafo de conflitos. Nas instâncias pequenas, em que o
ótimo é conhecido, a heurística de grau mínimo chega muito próximo dele.

% ---------------------------------------------------------------------------
\subsection{Defeitos em implementações públicas}
\label{sec:defeitos}

Um resultado que só emerge da execução repetida dos algoritmos --- e que um
levantamento puramente bibliográfico não revelaria --- diz respeito à
confiabilidade de implementações públicas de algoritmos clássicos. Duas das
quatro implementações do método Húngaro avaliadas apresentam defeitos.

A primeira (\texttt{hunalgorithm}) falha de forma determinística em matrizes
densas a partir de um certo tamanho, por exceder um limite interno de iterações
previsto em seu próprio código; matrizes estruturadas, que produzem padrões mais
simples, não acionam o problema.

A segunda (\texttt{benchaplin}) é mais notável, por ser \textbf{não-determinística
entre execuções}. Sobre exatamente a mesma matriz de entrada, a implementação ora
retorna o custo ótimo correto, ora falha com um erro de tipo, ora ingressa em um
caminho de execução tão lento que não conclui em tempo praticável. A causa é a
dependência do defeito quanto à ordem de iteração de uma coleção de objetos, que a
linguagem aleatoriza a cada processo. Trata-se de três desfechos distintos para a
mesma entrada, todos observados. Além de ilustrar de forma vívida o valor de
\textit{implementar e medir} sobre apenas catalogar, esse comportamento tem uma
consequência metodológica registrada no capítulo anterior: a fixação da semente
de aleatorização foi necessária justamente para tornar os resultados
reprodutíveis.

Optou-se por \textbf{reportar} esses defeitos, e não corrigi-los: como o objeto de
estudo é o conjunto de implementações efetivamente disponíveis, corrigi-los
apagaria um resultado legítimo sobre a qualidade do que se encontra publicado.

% ---------------------------------------------------------------------------
\subsection{Crescimento empírico frente à complexidade teórica}
\label{sec:crescimento}

Como o tempo absoluto não é comparável entre linguagens, a análise de crescimento
apoia-se no expoente empírico definido na Seção~\ref{sec:criterios}: a inclinação
da reta ajustada ao gráfico do tempo pelo tamanho, em escala logarítmica. As
Figuras~\ref{fig:cresc-mcm} e~\ref{fig:cresc-assignment} apresentam esses
gráficos para dois problemas representativos, e as
Tabelas~\ref{tab:exp-mcm} a~\ref{tab:exp-general-weight} confrontam o expoente
medido com a complexidade teórica de cada algoritmo.

\begin{figure}[H]
    \centering
    \begin{tikzpicture}
        \begin{loglogaxis}[
            width=0.88\textwidth, height=7.2cm,
            xlabel={$n$}, ylabel={tempo mediano (ms)},
            legend pos=north west, legend cell align=left,
            legend style={font=\scriptsize}, grid=both, font=\footnotesize,
            log basis x=10, log basis y=10]
            \addplot+[mark=*] table[x=n,y=t,col sep=comma] {results/plot/mcm-ek.csv};
            \addlegendentry{Edmonds-Karp $O(V E^2)$}
            \addplot+[mark=square*] table[x=n,y=t,col sep=comma] {results/plot/mcm-aug.csv};
            \addlegendentry{Caminhos aumentantes $O(VE)$}
            \addplot+[mark=triangle*] table[x=n,y=t,col sep=comma] {results/plot/mcm-ks.csv};
            \addlegendentry{Karp-Sipser (heur.)}
            \addplot+[mark=diamond*] table[x=n,y=t,col sep=comma] {results/plot/mcm-hk.csv};
            \addlegendentry{Hopcroft-Karp $O(E\sqrt{V})$}
            \addplot+[mark=o] table[x=n,y=t,col sep=comma] {results/plot/mcm-greedy.csv};
            \addlegendentry{Gulosa (heur.)}
        \end{loglogaxis}
    \end{tikzpicture}
    \caption{Tempo de execução em escala log-log para o grafo denso (cardinalidade
    máxima bipartida). A inclinação de cada reta corresponde ao expoente empírico
    de crescimento; a de Edmonds-Karp, próxima de 3, destaca-se das demais.}
    \label{fig:cresc-mcm}
\end{figure}


\begin{figure}[H]
    \centering
    \adjustbox{max width=\textwidth}{%
    \begin{tikzpicture}
        \begin{loglogaxis}[
            width=0.88\textwidth, height=7.2cm,
            xlabel={$n$}, ylabel={tempo mediano (ms)},
            legend pos=north west, legend cell align=left,
            legend style={font=\scriptsize}, grid=both, font=\footnotesize,
            log basis x=10, log basis y=10]
            \addplot+[mark=*] table[x=n,y=t,col sep=comma] {results/plot/assignment-munkres.csv};
            \addlegendentry{Húngaro (munkres)}
            \addplot+[mark=square*] table[x=n,y=t,col sep=comma] {results/plot/assignment-mayorx.csv};
            \addlegendentry{Húngaro (mayorx)}
            \addplot+[mark=triangle*] table[x=n,y=t,col sep=comma] {results/plot/assignment-ssp.csv};
            \addlegendentry{Caminhos mín.\ sucessivos}
            \addplot+[mark=diamond*] table[x=n,y=t,col sep=comma] {results/plot/assignment-scipy.csv};
            \addlegendentry{SciPy (compilado)}
            \addplot+[mark=o] table[x=n,y=t,col sep=comma] {results/plot/assignment-jv.csv};
            \addlegendentry{Jonker-Volgenant (compilado)}
        \end{loglogaxis}
    \end{tikzpicture}%
    }
    \caption{Tempo de execução em escala log-log para a matriz densa (atribuição
    linear). As duas implementações de núcleo compilado situam-se bem abaixo das
    interpretadas, mas com inclinação semelhante.}
    \label{fig:cresc-assignment}
\end{figure}


Os expoentes medidos são, em geral, coerentes com a teoria. O caso mais nítido é
o do Edmonds-Karp, cujo expoente próximo de três acompanha o comportamento cúbico
previsto e o separa visivelmente dos demais na Figura~\ref{fig:cresc-mcm}. Nas
implementações de núcleo compilado, a reta situa-se muito abaixo das
interpretadas, mas com inclinação semelhante --- exatamente o efeito esperado: a
linguagem altera a constante (a posição da reta), não o expoente (a inclinação).

\begin{table}[H]
    \centering
    \caption{Expoente empírico de crescimento --- Grafo 3 — Denso   | COM perfeito | MCM = n}
    \label{tab:exp-mcm}
    \renewcommand{\arraystretch}{1.25}
    \footnotesize
    \begin{tabularx}{\textwidth}{@{} X l c c @{}}
        \toprule
        \textbf{Algoritmo} & \textbf{Compl.\ teórica} & \textbf{Expoente medido} & \textbf{$R^2$} \\
        \midrule
        Aug. Paths (wbchristerson) & $O(V \cdot E)$ & 2.61 & 0.917 \\
        Hopcroft-Karp (sofiat) & $O(E \cdot \sqrt{V})$ & 1.20 & 0.996 \\
        Edmonds-Karp (Maxflow-Algs) & $O(V \cdot E^{2})$ & 2.79 & 0.999 \\
        Hopcroft-Karp (NetworkX) & $O(E \cdot \sqrt{V})$ & 1.09 & 0.999 \\
        Bipartite Match. (igraph) & $O(E \cdot \sqrt{V})$ & 1.05 & 0.994 \\
        Gulosa & $O(E)$ & 1.00 & 1.000 \\
        Karp-Sipser & $O(V \cdot E)$ & 1.68 & 0.983 \\
        \bottomrule
    \end{tabularx}
\end{table}

\begin{table}[H]
    \centering
    \caption{Expoente empírico de crescimento --- Matriz 1 — Densa uniforme  | custo in [1,100]}
    \label{tab:exp-assignment}
    \renewcommand{\arraystretch}{1.25}
    \footnotesize
    \begin{tabularx}{\textwidth}{@{} X l c c @{}}
        \toprule
        \textbf{Algoritmo} & \textbf{Compl.\ teórica} & \textbf{Expoente medido} & \textbf{$R^2$} \\
        \midrule
        Húngaro (munkres) & $O(n^{3})$ & 2.36 & 0.999 \\
        Húngaro (benchaplin) & $O(n^{3})$ & 3.22 & 0.998 \\
        Húngaro (mayorx/KM) & $O(n^{3})$ & 1.39 & 0.980 \\
        SSP (flows) & $O(n^{4})*$ & 2.71 & 0.998 \\
        SciPy (linear\_sum\_assg) & $O(n^{3})$ & 1.86 & 0.994 \\
        Jonker-Volgenant (lap) & $O(n^{3})$ & 1.65 & 0.987 \\
        \bottomrule
    \end{tabularx}
\end{table}


Observa-se ainda que, nos tamanhos avaliados, o expoente medido frequentemente
fica \textit{abaixo} do previsto pela complexidade de pior caso, que é
conservadora: o pior caso raramente se materializa nas instâncias típicas. A
exceção reveladora é a variante orientada a objetos do Gale-Shapley
(Tabela~\ref{tab:exp-stable-marriage}), cujo expoente medido supera o teórico: o
sobrecusto de implementação a torna, na prática, super-quadrática, confirmando por
medição a ressalva já registrada sobre ela.

\begin{table}[H]
    \centering
    \caption{Expoente empírico de crescimento --- Perfil 1 — Aleatório       | permutações uniformes}
    \label{tab:exp-stable-marriage}
    \renewcommand{\arraystretch}{1.25}
    \footnotesize
    \begin{tabularx}{\textwidth}{@{} X l c c @{}}
        \toprule
        \textbf{Algoritmo} & \textbf{Compl.\ teórica} & \textbf{Expoente medido} & \textbf{$R^2$} \\
        \midrule
        Gale-Shapley (OOP/pip) & $O(n^{2})†$ & 3.28 & 0.996 \\
        Gale-Shapley (numérico/pip) & $O(n^{2})$ & 1.94 & 0.999 \\
        Gale-Shapley (Lattas/dict) & $O(n^{3})*$ & 2.07 & 0.999 \\
        \bottomrule
    \end{tabularx}
\end{table}

\begin{table}[H]
    \centering
    \caption{Expoente empírico de crescimento --- Matriz 1 — Densa uniforme  | custo in [1,100]}
    \label{tab:exp-bottleneck}
    \renewcommand{\arraystretch}{1.25}
    \footnotesize
    \begin{tabularx}{\textwidth}{@{} X l c c @{}}
        \toprule
        \textbf{Algoritmo} & \textbf{Compl.\ teórica} & \textbf{Expoente medido} & \textbf{$R^2$} \\
        \midrule
        Threshold (busca binária) & $O(n^{3} log n)$ & 2.61 & 0.996 \\
        Caminhos aumentantes & $O(n^{3})$ & 2.41 & 0.999 \\
        \bottomrule
    \end{tabularx}
\end{table}

\begin{table}[H]
    \centering
    \caption{Expoente empírico de crescimento --- Peso — esparso (grau ~3)}
    \label{tab:exp-general-weight}
    \renewcommand{\arraystretch}{1.25}
    \footnotesize
    \begin{tabularx}{\textwidth}{@{} X l c c @{}}
        \toprule
        \textbf{Algoritmo} & \textbf{Compl.\ teórica} & \textbf{Expoente medido} & \textbf{$R^2$} \\
        \midrule
        van Rantwijk (mwmatching) & $O(V^{3})$ & 1.40 & 0.989 \\
        NetworkX (peso máximo) & $O(V^{3})$ & 1.36 & 0.990 \\
        Path Growing & $O(E)$ & 0.88 & 0.998 \\
        \bottomrule
    \end{tabularx}
\end{table}

